\begin{ExerciseBlock}
\ExerciseHeading{08}{Mehrere Federn}{forces_springs_in_series_and_parallel}
Zur Befestigung einer Masse $m$ an einer Platte stehen 2 Federn mit unterschiedlichen Federkonstanten $D_1$ und $D_2$ zur Verfügung.

\begin{enumerate}
	\item Berechnen Sie die Federkonstante $D_{ges}$ einer einzigen Feder, die sich so verhält, wie die beiden in Reihe verbauten Federn. Zeichnen Sie in eine Skizze die an den Federn anfallenden Kräfte ein.
\begin{center}
	\CachedTikzFigure{content/resources/chapters/forces/tikz/springs_in_series.tex}
\end{center}

	\item Berechnen Sie die Federkonstante $D_{ges}$ einer einzigen Feder, die sich so verhält als wären die beiden Federn parallel verbaut. Zeichnen Sie in eine Skizze die an den Federn anfallenden Kräfte ein.
\begin{center}
	\CachedTikzFigure{content/resources/chapters/forces/tikz/springs_in_parallel.tex}
\end{center}

\end{enumerate}
\end{ExerciseBlock}
\UseExerciseSolution{forces_springs_in_series_and_parallel}
